\subsubsection{C 类：数据口径、隔离验证与决策解释并重}

C 类从字段、样本单位、时间范围、缺失异常和训练/验证边界写起。有效表达包括
“以……为一条观测，按主体/时间划分训练与验证集”“所有标准化和特征选择仅在
训练折拟合”“相对持续性/线性基线，模型在……指标上……”“该变化主要对应……
特征，但结论限于……范围”。评价排序需解释指标正向化、权重来源和排序稳定性；
预测需把误差转换为补货、时点或风险决策的影响。

需要避开的是全数据清洗后再切分、用随机切分检验未来预测、把相关性写成因果、
只报告最优模型以及用二维降维图证明分类性能。图表以数据质量、分布与相关、
验证协议、残差/校准、特征影响和决策结果为主。

